\topic{1}{Zasady dynamiki Newtona, pęd i jego zachowanie}{
W układzie inercjalnym ciało zachowuje stan spoczynku lub ruchu
jednostajnego prostoliniowego, jeżeli siła wypadkowa jest zerowa.
Druga zasada ma ogólną postać $\mathbf F_{\rm wyp}=d\mathbf p/dt$;
przy stałej masie redukuje się do $m\mathbf a$.
Trzecia zasada mówi, że siły wzajemnego oddziaływania dwóch ciał
mają przeciwne zwroty i tę samą wartość, ale działają na różne ciała.

Dla układu ciał sumujemy pędy: $\mathbf P=\sum_i\mathbf p_i$.
Siły wewnętrzne znoszą się parami w bilansie całego układu,
więc $d\mathbf P/dt=\mathbf F_{\rm zewn}$. Gdy zewnętrzny impuls
jest pomijalny, $\mathbf P$ pozostaje stały. Dotyczy to także
zderzeń niesprężystych, w których energia kinetyczna może maleć.
Przed zastosowaniem zasady trzeba jasno określić granice układu
i sprawdzić, czy impuls sił zewnętrznych jest rzeczywiście mały.}
\topic{2}{Układy inercjalne i nieinercjalne, siły bezwładności}{
Układ inercjalny to taki, w którym swobodne ciało porusza się
jednostajnie po prostej. Układ poruszający się względem niego
ze stałą prędkością także jest inercjalny. W układzie
przyspieszającym równanie Newtona wymaga dodania sił pozornych.
Dla układu przesuwającego się z przyspieszeniem $\mathbf A$
pojawia się siła $-m\mathbf A$.

W układzie obracającym się z prędkością kątową $\boldsymbol\Omega$
występują między innymi siła Coriolisa
$-2m\boldsymbol\Omega\times\mathbf v'$ i odśrodkowa
$-m\boldsymbol\Omega\times(\boldsymbol\Omega\times\mathbf r)$.
Jeżeli $\boldsymbol\Omega$ się zmienia, dochodzi siła Eulera.
Siły bezwładności nie wynikają z oddziaływania z innym ciałem;
są skutkiem wyboru współrzędnych. Przykładem jest odchylenie
toru długiego ruchu na obracającej się Ziemi.}
\newpage
\topic{3}{Praca, energia i siły zachowawcze}{
Praca siły na drodze $C$ jest całką
$W=\int_C\mathbf F\cdot d\mathbf r$.
Z twierdzenia o pracy i energii kinetycznej wynika
$W_{\rm wyp}=\Delta K$, gdzie $K=mv^2/2$ dla ruchu
nierelatywistycznego. Siła zachowawcza ma pracę niezależną
od drogi i można ją zapisać jako $\mathbf F=-\nabla U$.
Przykładami są grawitacja Newtonowska i idealna sprężyna.

Gdy działają tylko siły zachowawcze, $E=K+U$ jest stałe.
Jeśli działa tarcie, $\Delta(K+U)=W_{\rm niezach}$; utracona
energia mechaniczna przechodzi zwykle w energię wewnętrzną.
W zadaniu trzeba określić poziom zerowy energii potencjalnej
i rozróżnić zachowanie energii całkowitej od zachowania
energii mechanicznej. Moc chwilowa wynosi
$P=dW/dt=\mathbf F\cdot\mathbf v$.}
\topic{4}{Dynamika bryły sztywnej i moment pędu}{
Bryła sztywna zachowuje odległości między punktami.
Jej ruch można rozłożyć na translację środka masy oraz obrót.
Dla obrotu wokół ustalonej osi moment bezwładności jest
$I=\sum_i m_i r_{\perp i}^2$, a energia kinetyczna obrotu
$K_{\rm obr}=I\omega^2/2$. Moment siły to
$\boldsymbol\tau=\mathbf r\times\mathbf F$.

Równanie dynamiki obrotowej ma postać
$d\mathbf L/dt=\boldsymbol\tau_{\rm zewn}$.
Przy zerowym zewnętrznym momencie siły moment pędu
$\mathbf L$ jest zachowany. Dla obrotu wokół osi głównej
$L=I\omega$, więc zmniejszenie $I$ zwiększa $\omega$.
W ogólnym ruchu przestrzennym $\mathbf L$ i
$\boldsymbol\omega$ nie muszą być równoległe, ponieważ
moment bezwładności jest tensorem.}
\newpage
\topic{5}{Oscylator harmoniczny, tłumienie, wymuszenie i rezonans}{
Idealny oscylator spełnia $m\ddot x+kx=0$.
Rozwiązanie ma postać $x=A\cos(\omega_0t+\phi)$,
gdzie $\omega_0=\sqrt{k/m}$. Po dodaniu oporu lepkiego
$b\dot x$ otrzymujemy $m\ddot x+b\dot x+kx=0$.
W zależności od $b^2-4mk$ odpowiedź jest oscylacyjna
tłumiona, krytyczna albo aperiodyczna.

Przy wymuszeniu $F_0\cos\Omega t$ rozwiązanie zawiera
zanikający stan przejściowy i stan ustalony. Jego amplituda to
$A(\Omega)=F_0/\sqrt{(k-m\Omega^2)^2+b^2\Omega^2}$.
Dla słabego tłumienia maksimum występuje blisko
$\omega_0$; dokładne położenie maksimum wychylenia zależy
od $b$. W rezonansie energia dostarczana w kolejnych okresach
kompensuje straty. Należy odróżniać rezonans amplitudy
wychylenia od rezonansu mocy.}
\topic{6}{Grawitacja Newtona, prawa Keplera i energia potencjalna}{
Dwie masy punktowe przyciągają się siłą
$F=Gm_1m_2/r^2$ skierowaną wzdłuż prostej łączącej ciała.
Przy zerze energii w nieskończoności
$U(r)=-Gm_1m_2/r$. Ujemna energia całkowita w problemie
dwóch ciał odpowiada orbicie związanej.

Pierwsze prawo Keplera mówi o elipsie ze Słońcem
w jednym ognisku. Drugie, o równych polach zakreślanych
w równych czasach, wynika z zachowania momentu pędu.
Trzecie daje $T^2=4\pi^2a^3/[G(M+m)]$ dla półosi wielkiej
$a$ w problemie dwóch ciał. Prawo odwrotnych kwadratów
wyjaśnia te zależności przy założeniu izolowanego układu.
W pobliżu powierzchni Ziemi przy małych zmianach wysokości
można stosować przybliżenie $U\simeq mgh$.}
\newpage
\topic{7}{Statyka i dynamika płynów}{
W cieczy spoczywającej ciśnienie rośnie z głębokością:
$p=p_0+\rho gh$. Prawo Pascala opisuje przekazywanie
zmiany ciśnienia w cieczy zamkniętej. Prawo Archimedesa
daje siłę wyporu równą ciężarowi wypartego płynu,
$F_w=\rho_{\mathrm{plyn}}gV_{\rm wyparte}$.

Dla ustalonego przepływu nieściśliwego równanie ciągłości
ma postać $Av=\mathrm{const}$ w pojedynczym przewodzie.
Wzdłuż linii prądu płynu idealnego zachodzi
$p+\rho v^2/2+\rho gz=\mathrm{const}$.
Równanie Bernoulliego wynika z bilansu energii, lecz wymaga
braku istotnych strat lepkości i braku pomp między
porównywanymi punktami. W praktyce przepływ turbulentny
i tarcie wymagają dodatkowego członu strat.}
\topic{8}{Transformacje Galileusza i Lorentza}{
Transformacja Galileusza dla względnej prędkości $V$
ma postać $x'=x-Vt$, $t'=t$. Zakłada wspólny czas i działa
dobrze, gdy prędkości są małe względem $c$.
Transformacja Lorentza to
$x'=\gamma(x-Vt)$,
$t'=\gamma(t-Vx/c^2)$, gdzie
$\gamma=(1-V^2/c^2)^{-1/2}$.

Z niej wynikają dylatacja czasu własnego
$\Delta t=\gamma\Delta\tau$ i skrócenie długości
w kierunku ruchu $L=L_0/\gamma$. Prędkości wzdłuż
tej osi składają się według
$u'=(u-V)/(1-uV/c^2)$. Pęd relatywistyczny
to $\mathbf p=\gamma m\mathbf v$, energia całkowita
$E=\gamma mc^2$, a kinetyczna $K=(\gamma-1)mc^2$.
W granicy $v\ll c$ wracamy do $mv^2/2$.}
\newpage
\topic{9}{Gaz doskonały i jego przemiany}{
Model gazu doskonałego zakłada cząstki o zaniedbywalnej
objętości i oddziaływaniach poza zderzeniami.
Równanie stanu $pV=nRT=Nk_BT$ wiąże wielkości
makroskopowe. Temperatura jest związana ze średnią
energią ruchu cząstek, a ciśnienie wynika z przekazu
pędu ścianom.

W przemianie izotermicznej gazu doskonałego $pV$
jest stałe. W izochorycznej stała jest objętość,
w izobarycznej ciśnienie. Dla odwracalnej przemiany
adiabatycznej $pV^\gamma=\mathrm{const}$, gdzie
$\gamma=C_p/C_V$; wzór nie opisuje dowolnego procesu
z izolacją cieplną, jeśli jest on silnie nieodwracalny.
Praca wykonana przez gaz wynosi
$W=\int p\,dV$, a pole pod wykresem $p(V)$
ma interpretację energetyczną.}
\topic{10}{Pierwsza zasada termodynamiki i ekwipartycja}{
Pierwsza zasada to bilans energii wewnętrznej.
Przy konwencji, że $W$ jest pracą wykonaną przez układ,
$\Delta U=Q-W$. Dla gazu doskonałego energia wewnętrzna
zależy tylko od temperatury; w przemianie izotermicznej
$\Delta U=0$. W przemianie adiabatycznej $Q=0$,
więc rozprężanie kosztem pracy obniża $U$.

Klasyczna zasada ekwipartycji przyznaje średnio
$k_BT/2$ każdemu niezależnemu kwadratowemu składnikowi
energii w równowadze. Dla gazu jednoatomowego
$U=3Nk_BT/2$ i $C_V=3Nk_B/2$.
Nie wolno jednak bezkrytycznie stosować jej do wszystkich
stopni swobody przy niskich temperaturach: wzbudzenia
rotacyjne i drganiowe mogą być zamrożone kwantowo.
Rozróżniaj pojemność cieplną całego układu od molowej.}
\newpage
\topic{11}{Druga zasada, ciepło zredukowane i entropia}{
Druga zasada wyznacza kierunek przemian: w izolowanym
układzie całkowita entropia nie maleje.
Dla odwracalnej drogi $\Delta S=\int\delta Q_{\rm rev}/T$;
iloraz $\delta Q/T$ nazywa się ciepłem zredukowanym.
Entropia $S$ jest funkcją stanu, choć samo ciepło
zależy od drogi.

Nierówność Clausiusa
$\oint\delta Q/T\leq0$ odróżnia proces odwracalny
od nieodwracalnego. W samorzutnym przepływie ciepła
od ciała ciepłego do zimnego suma zmian entropii
jest dodatnia. Maszyna cieplna nie może zamienić
całego ciepła pobranego z jednego zbiornika na pracę
w cyklu. Przy obliczeniach entropii procesu
nieodwracalnego wybieramy dogodną hipotetyczną drogę
odwracalną między tymi samymi stanami.}
\topic{12}{Napięcie powierzchniowe, włoskowatość i ruchy Browna}{
Siły spójności między cząsteczkami powodują, że
utworzenie nowej powierzchni cieczy wymaga energii.
Napięcie powierzchniowe $\sigma$ jest energią na
jednostkę powierzchni albo siłą na jednostkę długości
w odpowiedniej geometrii. Dla kulistej kropli różnica
ciśnień wynosi $\Delta p=2\sigma/R$.

W cienkiej rurce kapilarnej wysokość podciągania
jest w przybliżeniu
$h=2\sigma\cos\theta/(\rho gr)$.
Znak $\cos\theta$ rozstrzyga, czy ciecz podnosi się,
czy opada. Ruchy Browna to losowe ruchy zawieszonych
drobnych cząstek pod wpływem uderzeń cząsteczek
płynu. Dla dyfuzji jednowymiarowej
$\langle x^2\rangle=2Dt$; jest to średnia
po wielu realizacjach, a nie tor pojedynczej cząstki.}
\newpage
\topic{13}{Transport w gazach: dyfuzja, ciepło, lepkość i prąd}{
Dyfuzja jest przepływem cząstek wywołanym gradientem
stężenia; pierwsze prawo Ficka zapisujemy
$\mathbf J=-D\nabla n$. Przewodnictwo cieplne
ma analogiczną postać $\mathbf q=-\kappa\nabla T$.
Lepkość przenosi pęd między warstwami gazu,
a dla prostego ścinania naprężenie jest proporcjonalne
do gradientu prędkości.

W kinetycznym obrazie współczynniki transportu zależą
od średniej drogi swobodnej i średniej szybkości cząstek.
Przewodnictwo elektryczne wymaga nośników ładunku:
$\mathbf J_e=\sigma_e\mathbf E$ w zakresie prawa Ohma.
Nie należy utożsamiać przewodnictwa elektrycznego
gazu z przewodnictwem cieplnym; pierwsze silnie zależy
od jonizacji. Relacje liniowe obowiązują blisko
równowagi i przy niezbyt dużych gradientach.}
\topic{14}{Potencjał, pole elektryczne, strumień i prawo Gaussa}{
Natężenie pola jest siłą na jednostkowy dodatni
ładunek próbny: $\mathbf E=\mathbf F/q$.
W elektrostatyce $\mathbf E=-\nabla\varphi$,
a różnica potencjałów odpowiada pracy pola
na jednostkę ładunku ze zmienionym znakiem.
Strumień przez powierzchnię to
$\Phi_E=\int_S\mathbf E\cdot d\mathbf S$.

Prawo Gaussa w próżni mówi
$\oint_S\mathbf E\cdot d\mathbf S
=Q_{\rm wew}/\varepsilon_0$.
Jest zawsze prawdziwe dla zamkniętej powierzchni,
ale praktycznie pozwala wyliczyć pole tylko przy
odpowiedniej symetrii, np. sferycznej, cylindrycznej
lub płaskiej. W postaci lokalnej
$\nabla\cdot\mathbf E=\rho/\varepsilon_0$.
Zewnętrzne ładunki mogą dawać pole na powierzchni,
lecz ich całkowity wkład do strumienia zamkniętego jest zerowy.}
\newpage
\topic{15}{Dielektryki, polaryzacja i zjawiska pokrewne}{
W dielektryku pole przesuwa ładunki związane lub
ustawia dipole. Wektor polaryzacji $\mathbf P$
jest momentem dipolowym na jednostkę objętości,
a $\mathbf D=\varepsilon_0\mathbf E+\mathbf P$.
Prawo Gaussa dla $\mathbf D$ obejmuje tylko ładunek
swobodny. W liniowym, izotropowym materiale
$\mathbf D=\varepsilon\mathbf E$.

Dla kondensatora płaskiego wypełnionego jednorodnym
dielektrykiem $C=\varepsilon A/d$, jeśli pomijamy
efekty brzegowe. Elektrostrykcja opisuje odkształcenie
wywołane polem; piezoelektryczność wiąże naprężenie
z polaryzacją w kryształach o odpowiedniej symetrii.
Ferroelektryki mają spontaniczną polaryzację
i pętlę histerezy. Trzeba odróżniać ich nieliniową
odpowiedź od prostego prawa $\mathbf P=\chi\varepsilon_0\mathbf E$.}
\topic{16}{Klasyczny model przewodnictwa elektrycznego}{
W modelu Drudego elektrony przewodnictwa są
nośnikami o średnim czasie relaksacji $\tau$.
W stałym polu ich średnia prędkość dryfu
spełnia $\mathbf v_d=-e\tau\mathbf E/m$.
Dla koncentracji $n$ prąd ma postać
$\mathbf J=-ne\mathbf v_d=\sigma\mathbf E$,
gdzie $\sigma=ne^2\tau/m$.

Model wyjaśnia prawo Ohma i jakościowo efekt Halla,
ale nie opisuje poprawnie wszystkich własności
metali, ponieważ pomija statystykę Fermiego
i strukturę pasmową. Czas relaksacji jest parametrem
efektywnym zależnym od rozpraszania na drganiach
sieci i defektach. Wartość dryfu jest mała
w porównaniu z typowymi mikroskopowymi szybkościami
elektronów; sygnał elektryczny rozchodzi się
znacznie szybciej niż pojedynczy nośnik.}
\newpage
\topic{17}{Prawo Ohma, SEM, Kirchhoff, praca i moc}{
Lokalne prawo Ohma to $\mathbf J=\sigma\mathbf E$.
Dla jednorodnego przewodnika o długości $L$
i przekroju $A$ mamy $R=\rho L/A$, a na jego
końcach $U=IR$. Siła elektromotoryczna źródła
jest pracą sił nieelektrostatycznych na jednostkę
ładunku w pełnym obiegu.

Pierwsze prawo Kirchhoffa wynika z zachowania
ładunku: suma prądów wpływających do węzła
równa się sumie wypływających. Drugie, w obwodzie
skupionym bez zmiennego strumienia magnetycznego,
sumuje napięcia w oczku do zera. Moc na oporniku
$P=UI=I^2R=U^2/R$, a energia wydzielona
w czasie $t$ przy stałej mocy wynosi $Pt$.
Przy zmiennym strumieniu trzeba uwzględnić
indukowaną SEM.}
\topic{18}{Prawo Ampère'a, siła przewodników i indukcja}{
W magnetostatyce prawo Ampère'a ma postać
$\oint\mathbf B\cdot d\boldsymbol\ell=\mu_0I_{\rm wew}$.
Pole długiego prostego przewodnika wynosi
$B=\mu_0I/(2\pi r)$. Dwa równoległe przewodniki
z prądami w tym samym kierunku przyciągają się;
siła na jednostkę długości to
$F/L=\mu_0I_1I_2/(2\pi d)$.

Prawo Faradaya mówi, że zmiana strumienia
magnetycznego indukuje siłę elektromotoryczną
$\mathcal E=-d\Phi_B/dt$. Znak minus wyraża
regułę Lenza: prąd indukowany przeciwdziała
zmianie strumienia. Strumień może się zmieniać
przez ruch obwodu, zmianę pola lub powierzchni.
Przy polu zależnym od czasu pełne prawo
Ampère'a wymaga dodatkowego członu Maxwella.}
\newpage
\topic{19}{Materia w polu magnetycznym}{
Magnetyzacja $\mathbf M$ to moment magnetyczny
na jednostkę objętości. Diamagnetyki mają
zwykle małą ujemną podatność: pole indukuje
moment przeciwny do pola. Paramagnetyki mają
dodatnią podatność, bo momenty atomowe
częściowo ustawiają się z polem.

Ferromagnetyki mogą mieć spontaniczne uporządkowanie
momentów poniżej temperatury Curie. Ich odpowiedź
jest nieliniowa i zależy od historii pola,
co daje histerezę oraz domeny magnetyczne.
W prostych materiałach liniowych
$\mathbf B=\mu_0(\mathbf H+\mathbf M)$
i $\mathbf M=\chi_m\mathbf H$, lecz nie należy
stosować stałej $\chi_m$ do całej pętli ferromagnetyka.
Mikroskopowe pochodzenie ferromagnetyzmu wymaga
kwantowej energii wymiany, której klasyczny
obraz małych magnesów nie wyjaśnia.}
\topic{20}{Drgania w obwodach LC i RLC}{
W idealnym obwodzie LC ładunek kondensatora
spełnia $L\ddot q+q/C=0$ i oscyluje z częstością
$\omega_0=1/\sqrt{LC}$. Energia przechodzi
między polem elektrycznym kondensatora
$q^2/(2C)$ a polem magnetycznym cewki
$Li^2/2$. W szeregowym RLC dochodzi tłumienie:
$L\ddot q+R\dot q+q/C=0$.

Przy sinusoidalnym wymuszeniu prądu zmiennego
moduł impedancji szeregowego obwodu wynosi
$|Z|=\sqrt{R^2+(\omega L-1/\omega C)^2}$.
Maksymalny prąd występuje przy
$\omega=1/\sqrt{LC}$ w tym idealizowanym modelu.
Dobroć $Q=\omega_0L/R$ opisuje ostrość rezonansu
przy małym tłumieniu. Trzeba podać, czy mówimy
o obwodzie szeregowym czy równoległym.}
\newpage
\topic{21}{Superpozycja fal, fale stojące i dudnienia}{
W ośrodku liniowym wychylenia fal sumują się.
Dwie fale o tej samej częstości i zgodnych fazach
mogą dać wzmocnienie, a o fazach przeciwnych
wygaszenie. Dwie przeciwbieżne fale o tej samej
amplitudzie tworzą falę stojącą,
$y=2A\sin(kx)\cos(\omega t)$.
Węzły pozostają nieruchome, a strzałki drgają
z największą amplitudą.

Przy bliskich częstościach $f_1$ i $f_2$
amplituda wypadkowa zmienia się wolno,
a częstość dudnień wynosi $|f_1-f_2|$.
W rurach i strunach warunki brzegowe wybierają
dozwolone długości fal i częstości własne.
Superpozycja zakłada liniowość równania falowego;
przy dużych amplitudach w ośrodku nieliniowym
proste dodawanie rozwiązań może zawieść.}
\topic{22}{Interferencja, doświadczenie Younga i Michelsona}{
Dla dwóch spójnych źródeł różnica dróg $\Delta$
określa różnicę faz $2\pi\Delta/\lambda$.
Wzmocnienie zachodzi przy $\Delta=m\lambda$,
a wygaszenie przy $\Delta=(m+1/2)\lambda$
dla źródeł zgodnych w fazie. W dalekim polu
dwu szczelin o odległości $d$ mamy
$\Delta\simeq d\sin\theta$.

W doświadczeniu Younga rozstaw prążków
na ekranie odległym o $L$ jest przy małych
kątach $\Delta x\simeq\lambda L/d$.
Natężenie dwóch jednakowych fal wynosi
$I=4I_0\cos^2(\Delta\phi/2)$.
Interferometr Michelsona rozdziela wiązkę
na dwa ramiona i łączy ją ponownie.
Przesunięcie zwierciadła o $\delta$ zmienia
drogę optyczną o $2\delta$, co umożliwia
precyzyjne pomiary długości i zmian refrakcji.}
\newpage
\topic{23}{Efekt Dopplera i fala uderzeniowa}{
Efekt Dopplera jest zmianą obserwowanej
częstotliwości wskutek ruchu źródła lub odbiorcy.
Dla dźwięku w nieruchomym ośrodku
$f'=f(v+v_o)/(v-v_s)$ przy odpowiednio
ustalonych dodatnich kierunkach zbliżania.
Wzór zależy od prędkości względem ośrodka.

Dla światła w próżni obowiązuje relatywistyczny
efekt Dopplera; przy oddalaniu wzdłuż linii
widzenia $f'=f\sqrt{(1-\beta)/(1+\beta)}$.
Nie wymaga on materialnego ośrodka.
Gdy źródło dźwięku przekracza prędkość fali,
fronty tworzą stożek Macha z
$\sin\theta=v_{\rm fala}/v_{\mathrm{zrodla}}$.
Przechodząca przez obserwatora fala uderzeniowa
powoduje skok ciśnienia. Nie jest to po prostu
przypadek graniczny zwykłego wzoru na ton
odbierany od poddźwiękowego źródła.}
\topic{24}{Widmo fal elektromagnetycznych}{
Fale elektromagnetyczne w próżni mają
$c=\lambda f$ i przenoszą energię; foton ma
$E=hf$. Od małych do dużych częstości
wyróżnia się fale radiowe, mikrofale,
podczerwień, światło widzialne, ultrafiolet,
promieniowanie rentgenowskie i gamma.
Granice pasm są umowne i częściowo zachodzą
na siebie.

Różnice praktyczne wynikają z energii fotonów
i oddziaływań z materią. Mikrofale silnie
oddziałują z pewnymi dipolami i strukturami,
widzialne pobudza receptory oka, a promieniowanie
UV, X oraz gamma może jonizować materię,
zależnie od energii. Rentgenowskie i gamma
często rozróżnia się także według sposobu
powstania: pierwsze wiąże się z elektronami,
drugie z przemianami jądrowymi.
Wszystkie te pasma podlegają równaniom Maxwella.}
\newpage
